% Franka Scoop MPM scaling report — single-file Overleaf template
% Data: _scratch/reports/env_scaling/scoop_env_scaling_full_2026-06-11.csv (canonical merged table)
\documentclass[a4paper,11pt]{article}

\usepackage[margin=0.85in, top=1.3in, bottom=0.85in, headheight=18pt, headsep=14pt]{geometry}
\usepackage[T1]{fontenc}
\usepackage[utf8]{inputenc}
\usepackage[scaled=0.95]{helvet}
\usepackage{microtype}
\usepackage{xcolor}
\usepackage{booktabs}
\usepackage{array}
\usepackage{graphicx}
\usepackage{siunitx}
\sisetup{
  group-separator={,},
  group-minimum-digits=4,
  detect-weight=true,
  detect-family=true,
  table-number-alignment=center,
}
\usepackage{enumitem}
\usepackage{amsmath}
\usepackage{amssymb}
\usepackage{pgfplots}
\usepackage{tikz}
\usepackage{tcolorbox}
\tcbuselibrary{skins,breakable}
\usepackage{titlesec}
\usepackage{fancyhdr}
\usepackage{hyperref}

\renewcommand{\familydefault}{\sfdefault}
\pgfplotsset{compat=1.18}
\usetikzlibrary{calc}

% NVIDIA / report palette (matches Isaac Lab ecosystem)
\definecolor{NvidiaGreen}{HTML}{76B900}
\definecolor{NvidiaGreenDark}{HTML}{558B00}
\definecolor{ReportText}{HTML}{1A1A1A}
\definecolor{ReportMuted}{HTML}{4D4D4D}
\definecolor{SummaryGrey}{HTML}{F2F2F2}
\definecolor{RuleGrey}{HTML}{D9D9D9}
\definecolor{TabBlue}{HTML}{1F77B4}
\definecolor{TabGreen}{HTML}{2CA02C}
\definecolor{TabOrange}{HTML}{FF7F0E}
\definecolor{TabRed}{HTML}{D62728}

\color{ReportText}
\hypersetup{colorlinks=true, linkcolor=NvidiaGreenDark, urlcolor=NvidiaGreenDark}

\setlength{\parindent}{0pt}
\setlength{\parskip}{0.45em}
\setlength{\headheight}{22pt}
\setlength{\headsep}{18pt}

% Full-width green banner (page 1)
\newcommand{\NvidiaTopBanner}{%
  \tikz[remember picture,overlay]{%
    \fill[NvidiaGreen] (current page.north west) rectangle ([yshift=-1.05cm]current page.north east);
    \node[anchor=west, text=white, font=\bfseries\small]
      at ([xshift=0.85in, yshift=-0.55cm]current page.north west) {NVIDIA | ISAAC LAB};
    \node[anchor=east, text=white, font=\small]
      at ([xshift=-0.85in, yshift=-0.55cm]current page.north east) {Internal technical report};
  }%
}

% Slim green banner for running pages
\newcommand{\ReportRunHeader}{%
  \begin{tcolorbox}[
    enhanced,
    colback=NvidiaGreen,
    colframe=NvidiaGreen,
    boxrule=0pt,
    arc=0pt,
    left=6pt,
    right=6pt,
    top=4pt,
    bottom=4pt,
    width=\headwidth,
    halign=left,
  ]
  {\color{white}\footnotesize\textbf{NVIDIA | ISAAC LAB}\hfill
   Scaling Report | Franka Scoop MPM}
  \end{tcolorbox}%
}

\newcommand{\ReportRule}{{\color{RuleGrey}\rule{\linewidth}{0.5pt}}}

\newcommand{\ReportCaption}[1]{{\small\color{ReportMuted} #1\par}}

% Centered report tables with consistent spacing
\newenvironment{reporttable}{%
  \vspace{0.35em}%
  \begin{center}%
  \small
  \renewcommand{\arraystretch}{1.3}%
  \setlength{\tabcolsep}{11pt}%
}{%
  \end{center}%
  \vspace{0.55em}%
}

\newcommand{\tableheading}[1]{{\color{ReportMuted}#1\par\vspace{0.35em}}}

\newcommand{\ReportHeading}[1]{%
  \vspace{0.6em}%
  {\color{NvidiaGreen}\bfseries\large #1\par}%
  \vspace{0.25em}\ReportRule
  \vspace{0.45em}%
}

\newtcolorbox{executivesummarybox}{
  enhanced,
  breakable,
  colback=SummaryGrey,
  colframe=SummaryGrey,
  boxrule=0pt,
  borderline west={3.5pt}{0pt}{NvidiaGreen},
  left=10pt,
  right=10pt,
  top=6pt,
  bottom=6pt,
  arc=0pt,
  before skip=0.3em,
  after skip=0.55em,
}

\titleformat{\section}
  {\color{NvidiaGreen}\bfseries\large}
  {\thesection\ }
  {0pt}
  {}
  [\vspace{0.2em}\ReportRule\vspace{0.35em}]

\titleformat{\subsection}
  {\bfseries\normalsize\color{ReportText}}
  {\thesubsection\ }
  {0pt}
  {}

\titlespacing*{\section}{0pt}{0.9em}{0.5em}
\titlespacing*{\subsection}{0pt}{0.65em}{0.3em}

\pagestyle{fancy}
\fancyhf{}
\fancyhead[C]{\ReportRunHeader}
\fancyfoot[C]{\small\thepage}
\renewcommand{\headrulewidth}{0pt}

\fancypagestyle{firstpage}{%
  \fancyhf{}
  \fancyfoot[C]{\small\thepage}
  \renewcommand{\headrulewidth}{0pt}
}

\pgfplotsset{
  reportplot/.style={
    width=0.88\linewidth,
    height=6.4cm,
    tick label style={/pgf/number format/1000 sep={}},
    label style={font=\small},
    title style={font=\normalsize\bfseries},
    grid=both,
    grid style={gray!25},
    major grid style={gray!25},
    minor tick num=1,
    minor grid style={gray!12},
    legend style={font=\small, fill=white, draw=gray!40},
    enlarge x limits=0.03,
    enlarge y limits=0.06,
  },
  logx base2/.style={
    xmode=log,
    log basis x={2},
    xtick pos=left,
    xlabel style={font=\small},
  },
  scoop envaxis/.style={
    logx base2,
    xmin=0.8, xmax=1500,
    xtick={1,32,128,398,1200},
    xticklabels={1,32,128,398,1200},
  },
}


% Bold a number inside an siunitx S column
\newcommand{\B}{\bfseries}

\begin{document}
\thispagestyle{firstpage}
\NvidiaTopBanner

\vspace*{0.05cm}
{\small\color{ReportMuted}%
June 11, 2026 \quad|\quad Branch \texttt{max/newton-coupling-manager} (\texttt{add2551326})
\quad|\quad Prepared by Maximilian Krause\par}

\vspace{0.55em}
{\fontsize{20}{24}\selectfont\bfseries MPM Environment Throughput Scaling\par}
\vspace{0.2em}
{\normalsize\color{ReportMuted} Isaac-Scoop-Franka-v0 --- grid type $\times$ voxel size $\times$ parallel environments\par}

\vspace{0.4em}
\ReportRule
\vspace{0.45em}

\ReportHeading{Executive summary}

\begin{executivesummarybox}
\begin{itemize}[leftmargin=1.25em, itemsep=0.15em, topsep=0pt, parsep=0pt]
  \item \textbf{Use sparse at training scale.} At 15\,mm it reaches \textbf{965 env-steps/s @
        398 envs} with 8.5\,GiB steady memory; the fixed grid reaches 748 env-steps/s at
        28.6\,GiB.
  \item \textbf{Keep fixed for few-env debugging and the 20\,mm reference.} Its CUDA-graph floor
        wins at the small env counts measured, and 20\,mm fixed gives the highest raw rate measured
        (1{,}117 env-steps/s @ 598 envs).
  \item \textbf{Memory is the limiter.} Fixed-grid memory tracks particles at
        $\approx$\,10\,GiB/M, versus $\approx$\,2.5\,GiB/M for sparse; on this RTX 5090 that moves
        the practical ceiling from $\approx$\,2.7\,M particles to $\gtrsim$\,8.5\,M.
\end{itemize}
\end{executivesummarybox}

\section{Methodology}

Each point is one headless run on an idle NVIDIA GeForce RTX 5090 (32{,}607\,MiB VRAM, driver
595.71.05), launched through \texttt{scoop\_run.sh} with 20 warm-up steps and 120 measured steps.
The workload is the full zero-action RL environment step: four MPM substeps, MuJoCo arm solve,
Newton IK tracking, and particle-count reward/observation code. The repeated point
(20\,mm fixed @ 128 envs) reproduced within 0.5\%.

\textbf{Fixed} uses \texttt{grid\_padding=8}, CUDA graph replay, and
$\max(400\,\text{k},\, 0.35 \times \text{particles})$ active cells. \textbf{Sparse} runs eagerly
with \texttt{grid\_padding=0}. Particle counts are tied to voxel size: 2{,}261/env at 20\,mm,
5{,}817/env at 15\,mm, and 21{,}237/env at 10\,mm. Reported metrics are step time,
env-steps/s, startup time, and GPU memory (fixed: peak; sparse: steady/peak).

\newpage

\section{Results at a glance}
\label{sec:results}

\tableheading{Throughput across the full matrix. Each cell: step time (ms) / env-steps per second.
``OOM'' = out of memory at solver init (\texttt{Warp CUDA error 2}); ``---'' = not measured.}
\begin{reporttable}
\setlength{\tabcolsep}{5.5pt}%
\footnotesize
\resizebox{\linewidth}{!}{%
\begin{tabular}{@{}r cc cc cc cc cc@{}}
\toprule
& \multicolumn{2}{c}{\textbf{20\,mm fixed}} & \multicolumn{2}{c}{\textbf{15\,mm fixed}}
& \multicolumn{2}{c}{\textbf{10\,mm fixed}} & \multicolumn{2}{c}{\textbf{15\,mm sparse}}
& \multicolumn{2}{c}{\textbf{10\,mm sparse}} \\
\cmidrule(lr){2-3}\cmidrule(lr){4-5}\cmidrule(lr){6-7}\cmidrule(lr){8-9}\cmidrule(lr){10-11}
\textbf{envs} & {ms} & {sps} & {ms} & {sps} & {ms} & {sps} & {ms} & {sps} & {ms} & {sps} \\
\midrule
1    & 16.2  & 62   & 17.2  & 58  & 18.8  & 53  & 83.8   & 12  & 82.3  & 12  \\
8    & 42.6  & 188  & 55.1  & 145 & 108.3 & 74  & 89.0   & 90  & 87.4  & 92  \\
32   & 122.4 & 262  & 139.1 & 230 & 157.7 & 203 & 92.7   & 345 & 102.5 & 312 \\
48   & ---   & ---  & ---   & --- & ---   & --- & 96.8   & 496 & ---   & --- \\
128  & 213.9 & 598  & 191.8 & 667 & 452.5 & 283 & 149.8  & 854 & 249.0 & \textbf{514} \\
398  & 373.4 & 1066 & 532.1 & 748 & \multicolumn{2}{c}{OOM} & 412.6 & \textbf{965} & 929.5 & 428 \\
598  & 535.3 & \textbf{1117} & \multicolumn{2}{c}{OOM} & --- & --- & 693.1  & 863 & --- & --- \\
800  & 742.2 & 1078 & ---   & --- & ---   & --- & 1065.5 & 751 & --- & --- \\
1024 & \multicolumn{2}{c}{OOM} & --- & --- & --- & --- & ---  & --- & --- & --- \\
1200 & ---   & ---  & ---   & --- & ---   & --- & 2012.4 & 596 & --- & --- \\
\bottomrule
\end{tabular}%
}
\end{reporttable}

\tableheading{GPU memory (GiB) and startup time (s). Fixed-grid memory is the peak ($\approx$
steady); sparse memory is steady\,/\,peak (the difference is the init transient, $\leq$\,27\%).}
\begin{reporttable}
\setlength{\tabcolsep}{5.5pt}%
\footnotesize
\resizebox{\linewidth}{!}{%
\begin{tabular}{@{}r ccc cc ccc cc@{}}
\toprule
& \multicolumn{5}{c}{\textbf{Memory (GiB)}} & \multicolumn{5}{c}{\textbf{Startup (s)}} \\
\cmidrule(lr){2-6}\cmidrule(lr){7-11}
\textbf{envs} & {20\,f} & {15\,f} & {10\,f} & {15\,sp} & {10\,sp}
              & {20\,f} & {15\,f} & {10\,f} & {15\,sp} & {10\,sp} \\
\midrule
1    & 1.8  & 1.9  & 2.0  & 1.8\,/\,1.8   & 1.7\,/\,1.7 & 7   & 7  & 8  & 8   & 9  \\
8    & 3.5  & 4.4  & 7.5  & 2.0\,/\,2.0   & 2.1\,/\,2.2 & 8   & 8  & 10 & 8   & 10 \\
32   & 7.4  & 8.9  & 11.7 & 2.5\,/\,2.5   & 3.4\,/\,3.5 & 9   & 11 & 14 & 10  & 13 \\
48   & ---  & ---  & ---  & 2.7\,/\,2.8   & ---         & --- & --- & --- & 11 & --- \\
128  & 10.3 & 12.4 & 28.6 & 4.0\,/\,4.5   & 8.0\,/\,9.0 & 15  & 18 & 34 & 16  & 29 \\
398  & 15.3 & 28.6 & 31.1$^{\dagger}$ & 8.5\,/\,10.0 & 21.0\,/\,24.8 & 42 & 54 & --- & 50 & 95 \\
598  & 20.8 & 29.6$^{\dagger}$ & --- & 11.3\,/\,13.7 & --- & 80  & --- & --- & 96  & --- \\
800  & 27.5 & ---  & ---  & 14.2\,/\,17.8 & --- & 145 & --- & --- & 165 & --- \\
1024 & 31.2$^{\dagger}$ & --- & --- & ---  & --- & --- & --- & --- & --- & --- \\
1200 & ---  & ---  & ---  & 20.4\,/\,25.7 & --- & --- & --- & --- & 390 & --- \\
\bottomrule
\end{tabular}%
}
\end{reporttable}
\ReportCaption{$^{\dagger}$\,peak reached by the failed allocation attempt of an OOM run; the GPU
capacity is 31.8\,GiB. Columns: f $=$ fixed, sp $=$ sparse.}

\section{Throughput and step time}

\begin{center}
\begin{tikzpicture}
\begin{axis}[
  reportplot,
  scoop envaxis,
  title={Total throughput},
  xlabel={Number of parallel envs (log scale)},
  ylabel={Env-steps per second},
  ymin=0, ymax=1300,
  ytick={0,200,400,600,800,1000,1200},
  legend pos=south east,
  legend columns=2,
]
\addplot[TabGreen, thick, mark=*, mark size=1.8pt] coordinates {
  (1,62) (32,262) (128,598) (398,1066) (598,1117) (800,1078)
};
\addlegendentry{20\,mm fixed}
\addplot[TabBlue, thick, mark=square*, mark size=1.8pt] coordinates {
  (1,58) (32,230) (128,667) (398,748)
};
\addlegendentry{15\,mm fixed}
\addplot[TabOrange, thick, mark=triangle*, mark size=2.2pt] coordinates {
  (1,53) (32,203) (128,283)
};
\addlegendentry{10\,mm fixed}
\addplot[TabBlue, thick, dashed, mark=o, mark size=2.0pt] coordinates {
  (1,12) (32,345) (128,854) (398,965) (598,863) (800,751) (1200,596)
};
\addlegendentry{15\,mm sparse}
\addplot[TabOrange, thick, dashed, mark=triangle, mark size=2.4pt] coordinates {
  (1,12) (32,312) (128,514) (398,428)
};
\addlegendentry{10\,mm sparse}
\addplot[only marks, TabRed, mark=x, mark size=4pt, thick] coordinates {(398,330) (1024,1078)};
\node[font=\scriptsize, color=TabRed, anchor=south east] at (axis cs:380,345) {10\,mm fixed OOM};
\node[font=\scriptsize, color=TabRed, anchor=north east] at (axis cs:1024,1050) {20\,mm fixed OOM};
\end{axis}
\end{tikzpicture}

\ReportCaption{Throughput peaks: 20\,mm fixed at 1{,}117 env-steps/s, 15\,mm sparse at 965, and
10\,mm sparse at 514. Sparse overtakes fixed at the same voxel by 32 envs; red $\times$ markers
show OOM boundaries. The omitted 8- and 48-env points remain in Section~\ref{sec:results}.}
\end{center}

\begin{center}
\begin{tikzpicture}
\begin{axis}[
  reportplot,
  scoop envaxis,
  ymode=log, log basis y={10},
  title={Env step time},
  xlabel={Number of parallel envs (log scale)},
  ylabel={Env step time (ms, log scale)},
  ymin=10, ymax=2500,
  ytick={10,30,100,300,1000,2500},
  yticklabels={10,30,100,300,1000,2500},
  legend pos=south east,
  legend columns=2,
]
\addplot[TabGreen, thick, mark=*, mark size=1.8pt] coordinates {
  (1,16.2) (32,122.4) (128,213.9) (398,373.4) (598,535.3) (800,742.2)
};
\addlegendentry{20\,mm fixed}
\addplot[TabBlue, thick, mark=square*, mark size=1.8pt] coordinates {
  (1,17.2) (32,139.1) (128,191.8) (398,532.1)
};
\addlegendentry{15\,mm fixed}
\addplot[TabOrange, thick, mark=triangle*, mark size=2.2pt] coordinates {
  (1,18.8) (32,157.7) (128,452.5)
};
\addlegendentry{10\,mm fixed}
\addplot[TabBlue, thick, dashed, mark=o, mark size=2.0pt] coordinates {
  (1,83.8) (32,92.7) (128,149.8) (398,412.6) (598,693.1) (800,1065.5) (1200,2012.4)
};
\addlegendentry{15\,mm sparse}
\addplot[TabOrange, thick, dashed, mark=triangle, mark size=2.4pt] coordinates {
  (1,82.3) (32,102.5) (128,249.0) (398,929.5)
};
\addlegendentry{10\,mm sparse}
\end{axis}
\end{tikzpicture}

\ReportCaption{Fixed has the lower one-env floor ($\sim$16--19\,ms) because it can replay a CUDA
graph; sparse starts near $\sim$83\,ms but scales better once launch overhead is amortized.}
\end{center}

\section{Fixed vs.\ sparse}

\tableheading{Head-to-head at the 15\,mm training voxel. Bold marks the better grid per env count.}
\begin{reporttable}
\setlength{\tabcolsep}{7pt}%
\begin{tabular}{@{}r S[table-format=3.1] S[table-format=3.0] S[table-format=3.1] S[table-format=3.0] c@{}}
\toprule
& \multicolumn{2}{c}{\textbf{Fixed (CUDA graph)}} & \multicolumn{2}{c}{\textbf{Sparse (eager)}} & \\
\cmidrule(lr){2-3}\cmidrule(lr){4-5}
\textbf{num\_envs} & {\textbf{Step (ms)}} & {\textbf{Steps/s}} & {\textbf{Step (ms)}} & {\textbf{Steps/s}} & \textbf{Memory (GiB, peak/steady)} \\
\midrule
1   & \B 17.2  & \B 58  & 83.8   & 12  & 1.9 vs.\ 1.8 \\
8   & \B 55.1  & \B 145 & 89.0   & 90  & 4.4 vs.\ 2.0 \\
32  & 139.1 & 230 & \B 92.7   & \B 345 & 8.9 vs.\ 2.5 \\
128 & 191.8 & 667 & \B 149.8  & \B 854 & 12.4 vs.\ 4.0 \\
398 & 532.1 & 748 & \B 412.6  & \B 965 & 28.6 vs.\ 8.5 \\
\bottomrule
\end{tabular}
\end{reporttable}

At 15\,mm, fixed wins only at 1--8 envs; from 32 envs upward sparse is faster and much lighter on
memory. At 398 envs it is 29\% faster at $3.4\times$ less memory. The 10\,mm case is more severe:
sparse is already ahead by 8 envs and can run the 398-env, 8.45\,M-particle case that fixed cannot
initialize.

\section{GPU memory}

\begin{center}
\begin{tikzpicture}
\begin{axis}[
  reportplot,
  scoop envaxis,
  title={GPU memory},
  xlabel={Number of parallel envs (log scale)},
  ylabel={GPU memory (GiB)},
  ymin=0, ymax=35,
  ytick={0,5,10,15,20,25,30,35},
  legend style={at={(0.98,0.03)}, anchor=south east, font=\footnotesize,
    fill opacity=0.85, text opacity=1},
  legend columns=2,
]
\addplot[TabGreen, thick, mark=*, mark size=1.8pt] coordinates {
  (1,1.8) (32,7.4) (128,10.3) (398,15.3) (598,20.8) (800,27.5)
};
\addlegendentry{20\,mm fixed}
\addplot[TabBlue, thick, mark=square*, mark size=1.8pt] coordinates {
  (1,1.9) (32,8.9) (128,12.4) (398,28.6)
};
\addlegendentry{15\,mm fixed}
\addplot[TabOrange, thick, mark=triangle*, mark size=2.2pt] coordinates {
  (1,2.0) (32,11.7) (128,28.6)
};
\addlegendentry{10\,mm fixed}
\addplot[TabBlue, thick, dashed, mark=o, mark size=2.0pt] coordinates {
  (1,1.8) (32,2.5) (128,4.0) (398,8.5) (598,11.3) (800,14.2) (1200,20.4)
};
\addlegendentry{15\,mm sparse}
\addplot[TabOrange, thick, dashed, mark=triangle, mark size=2.4pt] coordinates {
  (1,1.7) (32,3.4) (128,8.0) (398,21.0)
};
\addlegendentry{10\,mm sparse}
\addplot[only marks, TabRed, mark=x, mark size=4pt, thick] coordinates {(398,31.1) (598,29.6) (1024,31.2)};
\draw[dashed, TabRed] (axis cs:0.8,31.8) -- (axis cs:1500,31.8);
\node[font=\scriptsize, color=TabRed, anchor=south east] at (axis cs:1450,32.2) {VRAM capacity (31.8\,GiB)};
\end{axis}
\end{tikzpicture}

\ReportCaption{Fixed-grid memory (solid; peak $\approx$ steady, includes the cell preallocation)
tracks total particles at $\approx$\,10\,GiB per million (least-squares: 4.7\,GiB $+$
10.4\,GiB/M); sparse steady memory (dashed) needs only $\approx$\,2.5\,GiB per million. Red
$\times$: failed allocation attempts of the OOM runs against the capacity line.}
\end{center}

Fixed memory follows total particles, not env count: 15\,mm @ 398 envs and 10\,mm @ 128 envs both
peak at 28.6\,GiB. Sparse pushes the same GPU beyond 8.5\,M particles and still leaves training
headroom for PPO and logging.

\section{Startup cost}

Startup is mostly flat through 48 envs ($\leq$\,16\,s), then scales with env count: 29--34\,s at
128 envs, 42--95\,s at 398, and 390\,s for sparse at 1{,}200. Even the largest case amortizes
within $\sim$200 steady-state env steps.

\section{Recommended operating points}

\begin{reporttable}
\begin{tabular}{@{}l c p{0.44\linewidth}@{}}
\toprule
\textbf{Goal} & \textbf{Config} & \textbf{Why} \\
\midrule
Training default        & 15\,mm sparse @ 398 envs & Fastest 15\,mm config: 965 env-steps/s at
                                              8.5\,GiB steady. \\
Balanced / shared GPU   & 15\,mm sparse @ 128 envs & 854 env-steps/s at 4.0\,GiB steady;
                                              faster than the best 15\,mm fixed run with far less
                                              memory. \\
Fine media (physics)    & 10\,mm sparse @ 128 envs & 514 env-steps/s at 8.0\,GiB steady; use
                                              398 envs only when particle count matters more than
                                              sample rate. \\
Few-env debugging       & 15\,mm fixed @ $\leq$8 envs & The CUDA-graph replay floor
                                              beats sparse at these counts. \\
Max raw sample rate     & 20\,mm fixed @ 598 envs & 1{,}117 env-steps/s --- the overall peak, but
                                              coarse for real scooping. \\
\bottomrule
\end{tabular}
\end{reporttable}

\section{Caveats}
\label{sec:caveats}

\begin{itemize}[leftmargin=1.4em, itemsep=0.18em]
  \item \textbf{Single measurement per point.} Each cell is one 120-step run. The repeated point
        reproduced within 0.5\%; treat $\lesssim$\,5\% gaps cautiously.
  \item \textbf{Zero-action workload.} The cup barely moves, so sparse sees near-best-case
        allocation behavior. A trained policy should raise sparse step times somewhat, but the
        scale ranking should hold.
  \item \textbf{Physics validity differs across voxels.} The hemisphere ladle wall is 24\,mm: at
        20\,mm it spans only 1.2 cells and will leak during real scooping. Treat 20\,mm as a
        performance reference, not a training setting.
  \item \textbf{Memory figures include the preallocation policy.} Fixed-grid
        \texttt{max\_active\_cells} was auto-scaled as
        $\max(400\text{k}, 0.35 \times \text{particles})$; a tighter cap saves memory but risks
        active-cell overflow.
\end{itemize}

\vspace{0.8em}
{\small\color{ReportMuted}%
Raw data: \texttt{\_scratch/reports/env\_scaling/scoop\_env\_scaling\_full\_2026-06-11.csv};
per-run JSON/stdout logs live beside it. Sweep drivers: \texttt{\_scratch/bench\_env\_scaling.py}
and \texttt{\_scratch/bench\_redo.py}; fragments are regenerated by
\texttt{\_scratch/reports/gen\_tex\_fragments.py}.}

\end{document}
